% TOP_PROBLEM: {"id": 7200017, "title": "Wikipedia geometry item 17: Section conjecture on splittings of group homomorphisms from fundamental groups of complete smo…", "category_id": 5, "category": {"id": 5, "name": "algebraic_geometry", "display_name": "Algebraic Geometry", "description": "Geometric objects defined by polynomial equations.", "slug": "algebraic-geometry", "order_index": 5, "created_at": "2026-07-31T15:26:25.671Z"}, "status": "open"}
% FIRST_POSTED: null
% ENTRY_KIND: statement_only
% Imported from: batch19.tex; UnsolvedMath ID 7200017
\documentclass[11pt]{article}
\usepackage{fontspec}
\setmainfont{Latin Modern Roman}
\setsansfont{Latin Modern Sans}
\usepackage{amsmath,amssymb,mathtools}
\usepackage{unicode-math}
\setmathfont{Latin Modern Math}
\usepackage{xurl}
\usepackage[hidelinks]{hyperref}
\newcommand{\sref}[2]{\hyperlink{source:#1:#2}{[S#2]}}
\newcommand{\cataloguescope}[1]{\par\noindent\textbf{Question.} #1\par}
\begin{document}
% UnsolvedMath ID 7200017; source catalogue code AMR-071-0017
\setcounter{section}{7200016}
\section{Wikipedia geometry item 17: Section conjecture on splittings of group homomorphisms from fundamental groups of complete smo…}\label{7200017}
{\small\sffamily UnsolvedMath ID 7200017 · Algebraic Geometry}
\subsection{Definitions and mathematical statement}

\paragraph{Shared notation.}$\mathbb N=\{1,2,\ldots\}$, and $\mathbb Z,\mathbb Q,\mathbb R,\mathbb C$ denote the integers, rationals, reals and complex numbers. Any other conventions are specified locally.


Let $k$ be a field finitely generated over $\mathbb Q$, fix an algebraic closure $\bar k$, and let $X/k$ be a smooth, proper, geometrically connected curve of genus at least two. Write $X_{\bar k}=X\times_k\bar k$, and choose a geometric base point $\bar x$. The étale fundamental groups fit into the exact sequence of profinite groups
\[1\longrightarrow\pi_1(X_{\bar k},\bar x)\longrightarrow\pi_1(X,\bar x)\mathrel{\mathop{\longrightarrow}^{p}}\operatorname{Gal}(\bar k/k)\longrightarrow1.\]
A section is a continuous homomorphism $s$ with $p\circ s=\mathrm{id}$. Two sections are equivalent if conjugate by an element of $\pi_1(X_{\bar k},\bar x)$. Every rational point $x\in X(k)$ induces a section class $[s_x]$, independently of the path used to compare base points.
\paragraph{Statement recorded in the supplied source.}
Is the map $X(k)\longrightarrow\{\text{continuous sections of }p\}/\pi_1(X_{\bar k},\bar x)$, $x\longmapsto[s_x]$, a bijection for every such $k$ and $X$? \sref{7200017}{2}
\subsection{Short English statement}
For a complete hyperbolic curve over a field finitely generated over the rationals, do its rational points correspond exactly to the splittings of its arithmetic fundamental-group sequence, up to conjugacy by the geometric fundamental group?
\subsection{Sources}
\begin{description}
\item[\hypertarget{source:7200017:1}{S1}] ULAM.AI and the UnsolvedMath Contributors, UnsolvedMath: A Curated Collection of Open Mathematics Problems, record 7200017 (AMR-071-0017). CC BY 4.0. \url{https://huggingface.co/datasets/ulamai/UnsolvedMath}.
\item[\hypertarget{source:7200017:2}{S2}] Hélène Esnault and Phùng Hô Hai, Packets in Grothendieck’s Section Conjecture (2007). Section 2, exact sequence (2.2) and the complete-curve formulation (SC), pp. 3–4. \url{https://arxiv.org/pdf/math/0703877v2}.
\end{description}
\label{end:7200017}
\end{document}
